\chapter{Generic constraint DSL}
\label{ch:dsl_generico_per_i_vincoli_en}

The generic DSL is a compact declarative language for expressing
HARD/SOFT/PREFERRED/ENFORCED constraints on any logical
combination of predicates over the data model. It is implemented
in \texttt{webui/backend/utils/general\_dsl.py} ($\sim$250 LOC,
recursive-descent parser + AST + post-hoc evaluator, no external
dependencies). The language is exposed through
\texttt{POST /api/constraints/general}, validated live by
\texttt{POST /api/constraints/general/validate}, and accessible
from the ``+ New DSL constraint'' modal of the Constraints tab.

\section{Philosophy}

Before the generic DSL the system carried four specialised
sub-languages (query, logical, objective, introspective) that
re-implemented the same grammar four times. The generic DSL
unifies all of them under a single grammar with four
\emph{flavours} of compilation (LIST, LOGIC, OBJECTIVE, EVAL):
one parser, many compilers. The syntax of
\texttt{forall x in S where Filter: Body} and of the atomic
predicates (\texttt{==}, \texttt{!=}, \texttt{<},
\texttt{in}, \dots) is identical across all contexts; only the
translation backend changes.

\section{BNF grammar}

\begin{lstlisting}[basicstyle=\ttfamily\small]
expr           ::= iff_expr
iff_expr       ::= implies_expr ( "<=>" implies_expr )*
implies_expr   ::= or_expr     ( "=>"  or_expr )*
or_expr        ::= and_expr    ( ("OR"|"or") and_expr )*
and_expr       ::= not_expr    ( ("AND"|"and") not_expr )*
not_expr       ::= ("NOT"|"not") not_expr | atom

atom           ::= quant_expr | count_expr | comparison
                 | call | bool_lit | "(" expr ")"

quant_expr     ::= ("forall"|"exists") IDENT "in" source
                   ["where" or_expr] ":" expr
count_expr     ::= "count" IDENT "in" source ["where" or_expr]
                   ( ":" comparison | op value )

source         ::= IDENT | IDENT ("." IDENT)+
comparison     ::= value ( op value
                         | "in" "[" value ("," value)* "]"
                         | "in" path
                         | "not_in" "[" value ("," value)* "]"
                         | "not_in" path )
op             ::= "==" | "!=" | "<" | "<=" | ">" | ">="
\end{lstlisting}

Iterable sources include \texttt{lessons}, \texttt{teachers},
\texttt{classes}, \texttt{classrooms}, \texttt{students},
\texttt{subjects}, \texttt{slots}, \texttt{days},
\texttt{hours}. Path-sources are also accepted (e.g.
\texttt{exists g in s.groups: g.name == "X"}).

\section{DSL $\to$ CP-SAT compilation (Step 1)}
\label{sec:dsl-cpsat-en}

Starting with Step 1 of the multi-day plan, the DSL ships a
\emph{compiler} into CP-SAT in
\texttt{engine/dsl\_to\_cpsat.py}, class
\texttt{DSLConstraintCompiler}. While the post-hoc evaluator
inspects an already-built solution to declare it feasible or
violated, the compiler adds constraints \emph{during} the
construction of the CP-SAT model so that the solver
structurally avoids violations.

The compiler is invoked at three sites:

\begin{itemize}
\item by \texttt{MonolithicSolver}
  (cap.~\ref{ch:metodo_cpsat_en}) when the
  \texttt{via\_dsl=True} flag is on;
\item by every Branch-and-Price pricer
  (cap.~\ref{ch:tecniche_avanzate_en}) so the same constraints
  reach the sub-problem;
\item by \texttt{PhaseBDaySolver} when called with
  \texttt{via\_dsl=True} to re-use the legacy pipeline
  augmented with DSL.
\end{itemize}

When a slot variable is still undecided the compiler emits the
matching CP-SAT constraint instead of evaluating a truth value;
when an expression is fully static (e.g. \texttt{consecutive(s1,
s2)} with both \texttt{s1, s2} bound by the enclosing
\texttt{forall}) the evaluation runs at compile time.

\section{Slot predicates: \texttt{consecutive}, \texttt{same\_day}}
\label{sec:dsl-slot-preds-en}

Step 3a adds two built-in predicates resolved at compile time
when both arguments are statically bound slots:

\begin{itemize}
\item \texttt{same\_day(s1, s2)} -- both slots fall on the same
  weekday.
\item \texttt{consecutive(s1, s2)} -- both slots fall on the
  same weekday and differ by exactly one hour ($|h_1-h_2|=1$).
\end{itemize}

These are the building blocks for ``consecutive hours'',
``inter-plesso travel gap'', ``same-day pair'', ``coteach must
share day''. Arguments accept either literal pairs
\texttt{(d, h)} or dotted paths like \texttt{l.slot} that
return a \texttt{(day, hour)} tuple.

\section{The \texttt{l.classroom.plesso} path}

Also in Step 3a, the post-hoc evaluator and the compiler
resolve the chained path \texttt{l.classroom.plesso}: per
lesson \texttt{l}, the assigned room is read from the
\texttt{classroom\_for\_slot} map provided by Phase B / the
classroom-assignment phase, and the plesso (campus) is then
resolved through the \texttt{plessi\_data} table. Without
those two structures the path returns \texttt{None} and the
compiler emits a diagnostic.

The canonical use is a commuting rule between campuses (see
the \texttt{vincoli} chapter, plessi section): ``no teacher
may have a lesson in Plesso B in the hour immediately
following a lesson in Plesso A''.

\section{Canonical pattern vocabulary (Step 3b)}
\label{sec:dsl-canonici-en}

Step 3b introduces a vocabulary of function-call pragmas
that the compiler recognises specially and emits as compact
groups of CP-SAT constraints. Each form corresponds to a
frequent pattern that, written as raw \texttt{forall} loops,
would produce hundreds or thousands of clauses. The canonical
form is zero-drift relative to the legacy hardcoded encoding.

\begin{tabular}{lp{6.7cm}}
\toprule
\textbf{Pragma} & \textbf{Legacy equivalent} \\
\midrule
\texttt{no\_holes\_class("3B")} &
  non-increasing chain forbidding gaps in 3B's weekly schedule
  (legacy \texttt{add\_class\_no\_holes}). \\
\texttt{class\_present\_at\_hour("3B", 11)} &
  if 3B has any lesson on a day, the slot at hour 11 is busy
  (legacy HARD-3, ``school day ends $\geq$ 12:00''). \\
\texttt{class\_day\_load\_in("3B", 0, 4, 5, 6)} &
  3B's daily load is 0 (free day) or 4--6 hours (legacy
  HARD-1 + HARD-2). \\
\texttt{teacher\_max\_per\_day("Rossi", 5)} &
  teacher Rossi has at most 5 hours on any day
  (\texttt{MAX\_PROF\_HOURS\_PER\_DAY}). \\
\texttt{cattedra\_max\_per\_day("Rossi", "3B", "Matematica", 2)} &
  the (Rossi, 3B, Math) cattedra has at most 2 hours per day
  (\texttt{MAX\_PER\_DAY\_TRIPLE}). \\
\texttt{subject\_pair\_must("3B", "Scienzemotorie")} &
  when 3B has 2 hours of PE on a day they MUST be consecutive
  (legacy HARD-B). \\
\texttt{subject\_pair\_exists("3B", "Matematica")} &
  when 3B has 2+ hours of math on a day at least one
  consecutive pair must exist (legacy HARD-A). \\
\bottomrule
\end{tabular}

\medskip

\texttt{no\_holes\_all\_classes()} and
\texttt{class\_present\_at\_hour\_all(11)} act in batch over
every class in the current slot scope.

\section{Plesso commuting via DSL (Step 3c)}

Step 3c wires the \texttt{PlessoCommutingRule} ORM table into
the DSL: the helper
\texttt{plesso\_commuting\_rule\_to\_dsl(rule)} returns a list
of DSL clauses equivalent to the legacy rule. A regression
test verifies zero-drift: applying the rule via DSL produces
exactly the same forbidden slot tuples as the hardcoded
\texttt{plessi\_constraints} pipeline.

\section{Seed legacy HARDs via DSL (Step 3d-3e)}

Step 3d-3e adds \texttt{seed\_implicit\_hardcoded(profs)}: given
the teacher dictionary, it returns the list of DSL strings that
encode every legacy HARD constraint of the solver. Combined
with the \texttt{via\_dsl=True} flag of \texttt{MonolithicSolver}
and \texttt{PhaseBDaySolver}, this allows building a complete
CP-SAT model \emph{exclusively from DSL} -- no hardcoded code
path is exercised. Invariants:

\begin{itemize}
\item \emph{syntactic zero-drift}: the slot tuples
  forbidden/forced by the DSL path coincide bit-for-bit with
  the legacy ones;
\item \emph{semantic zero-drift}: the feasible set accepted by
  the two paths is identical, validated on the
  \texttt{small} / \texttt{medium} / \texttt{huge} benchmarks;
\item \emph{determinism}: the order in which
  \texttt{seed\_implicit\_hardcoded} emits DSL strings is
  stable, so the CP-SAT model construction order is
  reproducible across runs.
\end{itemize}

This is the migration vehicle: in subsequent steps the
hardcoded path will be deprecated in favour of the DSL path,
and the existence of \texttt{seed\_implicit\_hardcoded}
allows that to happen without coverage loss.

\section{SOFT constraints with weight}
\label{sec:dsl-soft-status-en}

The DSL accepts four levels: \texttt{hard}, \texttt{soft},
\texttt{preferred}, \texttt{enforced}. The \emph{post-hoc
evaluator} fully implements SOFT: violations contribute to the
solution's global score and surface in \texttt{run\_metrics}.

Since the \texttt{feat(dsl): soft level for general\_dsl}
commit, the \emph{CP-SAT compiler} also wires SOFT directly
into the solver objective. Pipeline:

\begin{enumerate}
\item User picks \texttt{level=soft} and an integer weight $W$
  (e.g. 50) in the create modal.
\item \texttt{load\_all\_dsl\_constraints(db,
  include\_soft=True)} returns the rule with
  \texttt{is\_hard=False, weight=$W$}.
\item \texttt{add\_dsl\_constraint(expr, is\_hard=False,
  soft\_weight=$W$)} compiles the rule: each candidate
  \texttt{slot[k]} that would violate the body emits
  \texttt{(W, slot[k])} into \texttt{soft\_cost\_terms}
  instead of \texttt{slot[k] == 0}; in nested
  \texttt{forall}-\texttt{forall} bodies a reified
  \texttt{BoolVar b} with \texttt{b == s1 \&\& s2} is created
  and \texttt{(W, b)} is appended.
\item \texttt{MonolithicSolver.build()} adds
  \texttt{dsl\_soft\_cost\_terms} to \texttt{obj\_terms}
  before \texttt{Minimize(sum(obj\_terms))}.
\end{enumerate}

\textbf{Semantics}. A \emph{forall} violation costs $W$ per
lesson placed in a forbidden cell; a
\emph{forall}-\emph{forall} violation costs $W$ per
co-selected pair. Negative weights act as PREFERRED bonuses.

\textbf{Example}. ``Rossi should not have consecutive-day Math
in 3A'' (penalty 50 per pair):
\begin{lstlisting}
forall l1 in lessons where l1.teacher == "Rossi"
                       and l1.class == "3A":
    forall l2 in lessons where l2.teacher == "Rossi"
                           and l2.class == "3A":
        (l1.day == l2.day)
        or not consecutive_days(l1.day, l2.day)
\end{lstlisting}
Saved with \texttt{level=soft, weight=50, scope=global}: the
solver may schedule two hours on Tuesday + Wednesday paying
50, but if a Tuesday + Thursday solution exists it picks that
instead.

\section{Scope: where to save the constraint}
\label{sec:dsl-scope-en}

The ``+ New DSL constraint'' modal appears in two places and
yields rules with different \texttt{scope}. Rule of thumb:
\texttt{scope} is the rule's \emph{domain}; the \texttt{forall}
body is the \emph{filter} of application.

\begin{itemize}
\item \textbf{Constraints tab} (\texttt{/constraints}, modal
  with \texttt{scope="global"}, \texttt{owner\_id=null}).
  School-wide rule. Examples: ``no one teaches at 14:00'',
  ``no Saturday afternoons''.

\item \textbf{Teacher tab} (\texttt{/teachers/<id>}, ``Custom
  advanced DSL'' card, \texttt{scope="teacher"},
  \texttt{owner\_id=teacher\_id}). Per-teacher rule: the outer
  \texttt{forall} is \emph{typically} pre-filtered on
  \texttt{l.teacher == TeacherName}. Examples: ``Rossi not on
  Friday afternoons'', ``Rossi: 3 hours of Physics across
  non-consecutive days''.
\end{itemize}

In both cases the rule goes through
\texttt{POST /api/constraints/general}, is loaded by
\texttt{load\_all\_dsl\_constraints(db)}, and compiled by
\texttt{add\_all\_dsl\_constraints\_from\_db}; the
\texttt{scope\_kind} field is informational (it filters the
rules in Feasibility Check + reporting).

\subsection{Advanced DSL vs logical unavailability}

\texttt{LogicalUnavailability} (the ``Logical unavailability''
modal in the teacher tab) is a simpler sub-DSL designed for
unavailability windows (DNF over \texttt{(day, hour)} slots);
the system renders those into the generic DSL via
\texttt{logical\_unavailability\_to\_dsl}. The generic DSL is
strictly more expressive (\texttt{count}, \texttt{exists},
\texttt{forall}, \texttt{same\_day} /
\texttt{consecutive\_days}, arithmetic,
\texttt{l.classroom.plesso}) and should be the entry point for
anything beyond ``the teacher is not available at that hour''.

\section{Built-in predicates: \texttt{consecutive\_days},
\texttt{same\_day}, \texttt{consecutive}}
\label{sec:dsl-predicates-ref-en}

Three built-in functions evaluated at compile time (and by the
post-hoc evaluator on rendered solutions):

\begin{tabular}{lp{8cm}}
\toprule
\textbf{Predicate} & \textbf{Meaning} \\
\midrule
\texttt{same\_day(s1, s2)} &
  both slots fall on the same weekday. Arguments: tuples
  \texttt{(d, h)} or \texttt{l.slot}. \\

\texttt{consecutive(s1, s2)} &
  same weekday and adjacent \emph{hours} ($|h_1-h_2|=1$). Path:
  \texttt{l.slot}. \\

\texttt{consecutive\_days(d1, d2)} &
  adjacent \emph{days}, no wrap-around: true for
  $|d_1-d_2|=1$ with $d \in \{1,\dots,6\}$ (Mon..Sat); false
  for $\{$Sat, Mon$\}$. Arguments: \texttt{l.day} (int) or
  numeric / weekday-string literals. \\
\bottomrule
\end{tabular}

\section{Arithmetic in the DSL}
\label{sec:dsl-arithmetic-en}

The parser supports binary \texttt{+} and \texttt{-} on
integers: useful for index-dependent windows or cumulative
counts. Example: ``Rossi has the same hours in 3A on the first
3 days as in 3B on the last 3 days'':
\begin{lstlisting}
(count l in lessons where l.teacher == "Rossi"
   and l.class == "3A" and l.day <= 3: l)
== (count l in lessons where l.teacher == "Rossi"
      and l.class == "3B" and l.day >= 4: l)
\end{lstlisting}

Mixing booleans and numbers is rejected: the parser raises
\texttt{aritmetica non valida fra bool e numero} for typos
like \texttt{(l.hour == 9) + 1}.

\section{Real-world cases (Rossi and friends)}
\label{sec:dsl-rossi-en}

Every example below is covered by the integration tests in
\texttt{webui/backend/tests/test\_rossi\_general\_dsl\_integration.py}
and
\texttt{webui/backend/tests/test\_global\_dsl\_and\_soft\_and\_plesso\_commute.py}.

\subsection{Case (a) -- 3 hours of Physics in 3A on
non-consecutive days (HARD)}

\begin{lstlisting}
forall l1 in lessons where l1.teacher == "Rossi"
                       and l1.class == "3A"
                       and l1.subject == "Fisica":
    forall l2 in lessons where l2.teacher == "Rossi"
                           and l2.class == "3A"
                           and l2.subject == "Fisica":
        (l1.day == l2.day)
        or not consecutive_days(l1.day, l2.day)
\end{lstlisting}

Saved from the teacher tab with \texttt{scope="teacher",
owner\_id=Rossi.id}. The compiler iterates Rossi/3A/Fisica
candidate slots and, for each pair $(l_1,l_2)$ with
\texttt{consecutive\_days($l_1$.day, $l_2$.day)} and different
days, emits \texttt{BoolOr([slot[k1].Not(),
slot[k2].Not()])}: the two slots cannot be co-selected.

\subsection{Case (b) -- 3 hours of Physics in 3C all on the
same day (HARD)}

\begin{lstlisting}
forall l1 in lessons where l1.teacher == "Rossi"
                       and l1.class == "3C"
                       and l1.subject == "Fisica":
    forall l2 in lessons where l2.teacher == "Rossi"
                           and l2.class == "3C"
                           and l2.subject == "Fisica":
        same_day(l1.slot, l2.slot)
\end{lstlisting}

By transitivity of \texttt{same\_day} under
\texttt{forall forall}, every pair of Rossi/3C/Fisica lessons
must share a day; the solver picks one day and packs all 3
hours there.

\subsection{Plesso commute case -- Rossi never A@9 $\to$
B@10 same day (HARD)}

\begin{lstlisting}
forall l1 in lessons where l1.teacher == "Rossi"
                       and l1.hour == 9
                       and l1.classroom.plesso == 1:
    forall l2 in lessons where l2.teacher == "Rossi"
                           and l2.hour == 10
                           and l2.classroom.plesso == 2
                           and same_day(l1.slot, l2.slot):
        false
\end{lstlisting}

Plesso 1 = ``A'', plesso 2 = ``B'' (the integer ids of
\texttt{plessi.id}). Compile-time when the model is built with
\texttt{classroom\_for\_slot} (e.g. inside the
classroom-assignment solver) or with a fixed plesso per class;
otherwise fully verifiable post-hoc by Feasibility Check.

\subsection{Case (Rossi no Physics on Friday) (HARD)}

\begin{lstlisting}
forall l in lessons where l.teacher == "Rossi"
                       and l.subject == "Fisica"
                       and l.day == 5:
    false
\end{lstlisting}

The \emph{static} \texttt{false} body is the canonical idiom
for ``forbid every slot satisfying the where''. Equivalent to
\texttt{l.day != 5} but more explicit when the where is
multi-clause.

\subsection{Case (same-day support coteach) (HARD)}

``When Rossi teaches Sostegno to 2B, the Sostegno lesson must
fall on the same day as a Math lesson Rossi gives to 2B'':
\begin{lstlisting}
forall l in lessons where l.teacher == "Rossi"
                       and l.class == "2B"
                       and l.subject == "Sostegno":
    exists m in lessons where m.teacher == "Rossi"
                          and m.class == "2B"
                          and m.subject == "Matematica"
                          and same_day(l.slot, m.slot):
        true
\end{lstlisting}

\subsection{Every case above as SOFT}

Switch \texttt{level=hard} to \texttt{level=soft} in the modal
and pick \texttt{weight=$W$}. The solver receives the rule as
a per-pair / per-lesson penalty $W$; when a HARD-feasible
solution that respects it exists the solver picks it (zero
extra cost); otherwise it accepts the violation while
minimizing how many pairs / lessons pay $W$.

\section{DSL architecture diagram}

\begin{verbatim}
+--------------------------------------------------+
|  UI ("+ New DSL constraint" modal)                |
|  webui/frontend/.../GeneralConstraintsCard.svelte |
+--------------------+-----------------------------+
                     |
                     v
+--------------------+-----------------------------+
|  POST /api/constraints/general                   |
+--------------------+-----------------------------+
                     |
                     v
+--------------------+-----------------------------+
|  Parser + AST -- general_dsl.py                  |
+--------------------+-----------------------------+
                     |
                     +---------+---------+
                     v         v         v
              post-hoc    CP-SAT      SOFT score
              evaluator   compiler    (post-hoc)
              (HARD       (Step 1+)
               viol. capture)
                     |         |
                     v         v
              +--------------------+
              |   CP-SAT solver    |
              |   (Mono / Pricer / |
              |    PhaseB / RF /   |
              |    Meta)           |
              +--------------------+
\end{verbatim}

One AST, two backends: the evaluator reads a solution and
answers ``feasible / violated''; the compiler emits CP-SAT
constraints before solving. The grammar is unique and source
names are identical between backends.
